% Generated draft: requires mathematical and source review.
\documentclass[11pt]{article}
\usepackage[margin=25mm]{geometry}
\usepackage{fontspec}
\setmainfont{Latin Modern Roman}
\usepackage{amsmath,amssymb,mathtools,unicode-math}
\setmathfont{Latin Modern Math}
\usepackage{enumitem,needspace,xurl,hyperref}
\setlength{\parindent}{0pt}
\setlength{\parskip}{5pt}
\setlength{\emergencystretch}{3em}
\newcommand{\N}{\mathbb{N}}
\newcommand{\Z}{\mathbb{Z}}
\newcommand{\Q}{\mathbb{Q}}
\newcommand{\R}{\mathbb{R}}
\newcommand{\C}{\mathbb{C}}
\newcommand{\F}{\mathbb{F}}
\newcommand{\A}{\mathbb{A}}
\newcommand{\PP}{\mathbb{P}}
\newcommand{\E}{\mathbb{E}}
\newcommand{\eps}{\varepsilon}
\newcommand{\dd}{\,\mathrm{d}}
\begin{document}
\section*{UnsolvedMath ID: 30006133}
\textbf{Source problem code:} catalogue entry
\subsection{Definitions and mathematical statement}
Let \(\mathbb F\in\{\mathbb C,\mathbb H\}\), let \(p=\dim_{\mathbb R}\mathbb F\in\{2,4\}\), and let \(m,n\) be integers satisfying \(1\leq m<n\). Identify \(\mathbb F^n\) with its underlying Euclidean space \(\mathbb R^{np}\). A convex body is a compact convex subset with nonempty interior. Let \(K\subset\mathbb F^n\) be such a body, with no symmetry or centering assumption. Write \(|K|\) for its \(np\)-dimensional Lebesgue volume, and set
\[
\kappa_j=\bigl|\{x\in\mathbb R^j:\|x\|\leq1\}\bigr|.
\]
Let \(\operatorname{Gr}_m(n,\mathbb F)\) be the Grassmannian of \(m\)-dimensional \(\mathbb F\)-linear subspaces of \(\mathbb F^n\), equipped with its invariant probability measure \(dE\). For \(E\in\operatorname{Gr}_m(n,\mathbb F)\), let \(P_E\) denote Euclidean orthogonal projection onto \(E\), and let \(|P_EK|\) denote \(mp\)-dimensional Lebesgue volume within \(E\).

The supplied record asks whether, for every such \(\mathbb F,n,m,K\),
\begin{equation}\tag{A}
|K|^{-m}\geq
\frac{\kappa_{mp}^{\,n}}{\kappa_{np}^{\,m}}
\int_{\operatorname{Gr}_m(n,\mathbb F)}|P_EK|^{-n}\,dE.
\end{equation}
It further requests an equality characterization in (A), stated as equality exactly for ``real, complex, or quaternionic ellipsoids.'' This wording is retained as an unresolved qualification. An ordinary real ellipsoid is a translate of an invertible real-linear image of the Euclidean unit ball; a possible meaning of an \(\mathbb F\)-ellipsoid is a translate of an invertible \(\mathbb F\)-linear image of that ball. The record does not specify which class is intended for each permitted field, so it does not supply a precise equality characterization.

The associated projection inequality is supplied as
\begin{equation}\tag{B, supplied}
\frac{\kappa_{np}}{\kappa_{mp}}
\int_{\operatorname{Gr}_m(n,\mathbb F)}|P_EK|\,dE
\geq |K|^{m/n}.
\end{equation}
There is an unresolved normalization defect in this second target. For the Euclidean unit ball, its asserted inequality becomes \(\kappa_{np}\geq\kappa_{np}^{m/n}\), which fails when \(\kappa_{np}<1\), including the permitted case \(\mathbb F=\mathbb C\), \(n=7\), since \(\kappa_{14}=\pi^7/5040<1\). Thus (B, supplied) cannot hold for all the stated parameters. The factor \(\kappa_{np}^{m/n}/\kappa_{mp}\) would give the unit-ball normalization consistent with (A); this is a possible correction requiring source verification and is not substituted for the supplied formula. Accordingly, the requested pair of inequalities and its equality clause remain a source-qualified draft requiring clarification.
\subsection{Short English statement}
For every full-dimensional convex body in complex or quaternionic Euclidean space, the source asks for an inequality involving negative powers of projection volumes and an associated inequality for their average. The supplied average-projection formula has an inconsistent normalization, and the asserted ellipsoid equality class is unspecified. These issues require clarification before the complete target can be stated as a consistent conjecture.
\subsection{Sources}
\begin{itemize}
\item Local UnsolvedMath registry: supplied complete source record, permanent ID \(30006133\), title \emph{Complex and Quaternionic Affine Quermassintegral Inequalities}, exact source problem code \texttt{catalogue entry}, in the Oberwolfach Reports Open Problems collection. No individual proposer is supplied; the proposed year is recorded as 2024.
\item Supplied report citation: \emph{Convex Geometry and its Applications} (2025), Oberwolfach Reports, pp.\ 3301--3376, DOI \texttt{10.4171/owr/2024/58}, \url{https://doi.org/10.4171/owr/2024/58}. The supplied original excerpt describes the projection inequalities as consequences of a conjecture and says they remain open over the complex numbers and quaternions. Its mathematical formula is textually degraded.
\item Supplied literature reference: Saroglou--Wannerer, \emph{Math. Ann.} 2025, \url{https://doi.org/10.1007/s00208-025-03279-3}.
\item Literature assessment checked at \texttt{2026-08-22}: the supplied curation labels the problem open and states that the real affine-quermassintegral inequality is known while the complex and quaternionic analogues remain open. This is supplied dated curation derived from the Oberwolfach record, not an independently verified literature assessment. The record itself requests independent verification of time-sensitive status and cited claims.
\item The registry labels the statement \texttt{corrected\_verified} and reports that it was reconstructed and checked against the official report. These are supplied provenance claims; no source was independently visited or verified for this draft.
\end{itemize}
\end{document}
